\subsection{Terminología}

\subsubsection{Terminología del negocio}

\begin{enumerate}
  \item \textbf{CAETEC:} Campo Agropecuario Experimental del Tecnológico de Monterrey. Es el socio formador del proyecto y funciona como espacio de investigación, experimentación y aplicación de tecnologías en actividades agropecuarias.
  \item \textbf{Ganado lechero:} Conjunto de bovinos destinados principalmente a la producción de leche.
  \item \textbf{Vaca lechera:} Bovino hembra criado y manejado para la producción de leche.
  \item \textbf{Ordeño:} Proceso mediante el cual se extrae la leche de una vaca.
  \item \textbf{Sistema de ordeño automático:} Sistema tecnológico que permite realizar el ordeño con una intervención humana mínima, utilizando sensores, mecanismos y software para controlar el proceso.
  \item \textbf{Cola:} Acumulación de individuos que esperan para acceder a un recurso o servicio que tiene una capacidad limitada.
  \item \textbf{Producción de leche:} Cantidad de leche producida por una vaca o por el conjunto del ganado durante un periodo determinado.
  \item \textbf{Condición corporal (BCS):} Evaluación del nivel de reservas corporales de una vaca mediante características visuales y/o físicas hecha por algún experto. BCS significa Body Condition Score.
\end{enumerate}

\subsubsection{Terminología de ciencia de datos e inteligencia artificial}

\begin{enumerate}
  \item \textbf{Datos:} Información que puede ser recopilada, almacenada y analizada para obtener conocimiento.
  \item \textbf{Dataset:} Conjunto estructurado de datos utilizado para análisis o desarrollo de modelos.
  \item \textbf{Variable:} Característica o atributo que puede tomar diferentes valores dentro de un conjunto de datos.
  \item \textbf{Feature:} Variable utilizada como entrada para un modelo de aprendizaje automático.
  \item \textbf{Variable objetivo:} Variable que un modelo intenta predecir o estimar.
  \item \textbf{Minería de datos:} Proceso de explorar grandes cantidades de datos para identificar patrones, relaciones o información útil.
  \item \textbf{Ciencia de datos:} Disciplina que combina estadística, programación, matemáticas y conocimiento del dominio para extraer información útil de los datos.
  \item \textbf{Inteligencia artificial (IA):} Área de la computación dedicada al desarrollo de sistemas capaces de realizar tareas que normalmente requieren capacidades asociadas con la inteligencia humana.
  \item \textbf{Aprendizaje automático (Machine Learning):} Subcampo de la IA que permite que los sistemas aprendan patrones a partir de datos para realizar predicciones o tomar decisiones.
  \item \textbf{Aprendizaje profundo (Deep Learning):} Subconjunto del aprendizaje automático dentro de la inteligencia artificial, el cual imita el funcionamiento del cerebro humano mediante el uso de redes neuronales artificiales de múltiples capas.
  \item \textbf{Aprendizaje supervisado:} Tipo de aprendizaje automático en el que el modelo aprende utilizando datos que contienen una respuesta o etiqueta conocida.
  \item \textbf{Modelo:} Representación matemática o computacional de patrones aprendidos a partir de datos, utilizada para realizar predicciones o clasificaciones.
  \item \textbf{Entrenamiento:} Proceso mediante el cual un modelo aprende los patrones presentes en los datos.
  \item \textbf{Predicción:} Resultado generado por un modelo para un nuevo dato que no necesariamente fue utilizado durante el entrenamiento.
  \item \textbf{Clasificación:} Tarea de aprendizaje automático en la que se asigna una observación a una categoría.
  \item \textbf{Regresión:} Tarea en la que un modelo estima un valor numérico.
  \item \textbf{Red neuronal artificial:} Modelo computacional inspirado en el funcionamiento de las redes neuronales biológicas, compuesto por capas de unidades que transforman información.
  \item \textbf{CNN (Convolutional Neural Network):} Tipo de red neuronal especializada en procesar imágenes y detectar patrones espaciales como bordes, formas y estructuras.
  \item \textbf{Visión por computadora:} Área de la IA que permite a las computadoras procesar e interpretar información contenida en imágenes o videos.
  \item \textbf{Procesamiento de imágenes:} Conjunto de técnicas utilizadas para transformar, mejorar o preparar imágenes para su análisis.
  \item \textbf{Etiqueta (label):} Valor asociado a un dato que representa la respuesta conocida que se desea aprender o predecir.
  \item \textbf{Conjunto de entrenamiento:} Parte del dataset utilizada para entrenar el modelo.
  \item \textbf{Conjunto de validación:} Datos utilizados durante el desarrollo para evaluar y ajustar el modelo sin utilizarlos directamente para entrenarlo.
  \item \textbf{Conjunto de prueba:} Datos que se mantienen separados del entrenamiento para evaluar el desempeño final del modelo.
  \item \textbf{Overfitting:} Situación en la que un modelo aprende demasiado los datos de entrenamiento y tiene un desempeño considerablemente peor con datos nuevos.
  \item \textbf{Underfitting:} Situación en la que un modelo es demasiado simple para capturar los patrones de los datos.
  \item \textbf{Generalización:} Capacidad de un modelo para producir resultados adecuados sobre datos nuevos que no utilizó durante su entrenamiento.
\end{enumerate}

\subsubsection{Terminología estadística}

\begin{enumerate}
  \item \textbf{Estadística:} Disciplina que permite recopilar, organizar, analizar e interpretar datos para obtener conclusiones.
  \item \textbf{Población:} Conjunto total de individuos u observaciones que son objeto de estudio.
  \item \textbf{Muestra:} Subconjunto de una población utilizado para realizar un análisis.
  \item \textbf{Media:} Valor obtenido al sumar todas las observaciones y dividirlas entre su cantidad.
  \item \textbf{Mediana:} Valor que se encuentra en el centro de un conjunto de datos ordenado.
  \item \textbf{Desviación estándar:} Medida que indica qué tan dispersos están los datos respecto a su media.
  \item \textbf{Correlación:} Medida de la relación entre dos variables.
  \item \textbf{Distribución:} Forma en que se encuentran organizados los valores de una variable.
  \item \textbf{Métrica:} Medida utilizada para cuantificar el desempeño de un modelo o proceso.
  \item \textbf{MAE (Mean Absolute Error):} Promedio del valor absoluto de las diferencias entre las predicciones y los valores reales.
  \item \textbf{RMSE (Root Mean Squared Error):} Raíz cuadrada del promedio de los errores al cuadrado.
  \item \textbf{Matriz de confusión:} Tabla que permite comparar las clases predichas por un modelo con las clases reales.
  \item \textbf{Exactitud (accuracy):} Proporción de predicciones correctas respecto al total de predicciones.
  \item \textbf{Error:} Diferencia entre un valor real y el valor estimado por un modelo.
\end{enumerate}
